\section{The Invariant Potential and Vacuum Selection}
\label{sec:potential:main}
\label{sec:vacuum}

The order parameter of UCFT is a field $\Phi$ valued in the matter representation
$27=J_3(\mathbb O)$ of $E_6$ (Axiom~I). This section presents the $E_6$-invariant potential
$V$ of postulate~\textbf{P1} and proves the vacuum-selection statement~\textbf{P2}: the
global-minimum locus of $V$ is exactly the rank-$1$ primitive-idempotent orbit
$\mathcal M=\mathrm{EIII}=E_6/(\mathrm{Spin}(10)\times U(1))$ of real dimension $32$. We also
show that the symmetric origin $\Phi=0$ is a strict local \emph{maximum}, which supplies the
roll-down dynamics. The potential class is the added postulate; that the chosen member of the
class selects the rank-$1$ orbit is a \textbf{theorem}. Throughout, the dynamics live on the
\emph{compact} real form $E_6$ (equivalently $E_{6(-14)}$, maximal compact
$H=\mathrm{Spin}(10)\times U(1)$); the non-compact $E_{6(-26)}$ enters only as the rigid
orbit classifier (rank and $N$ only). The vacuum geometry and its tangent space are used in
\autoref{sec:spacetime:main} (signature) and \autoref{sec:gravity:main} (spectrum, induced
gravity).

\subsection{Albert-algebra setup and the two real forms}
\label{sec:potential:setup}

Let $\mathbb O$ be the real octonions with conjugation $x\mapsto\bar x$ and norm
$n(x)=x\bar x$. The real Albert algebra is the $27$-dimensional Euclidean (formally real)
Jordan algebra
\begin{equation}
\label{eq:potential:albert}
J_3(\mathbb O)=\left\{
X=\begin{pmatrix}\xi_1 & x_3 & \bar x_2\\ \bar x_3 & \xi_2 & x_1\\ x_2 & \bar x_1 & \xi_3\end{pmatrix}
:\ \xi_i\in\mathbb R,\ x_i\in\mathbb O\right\},
\qquad X\circ Y=\tfrac12(XY+YX).
\end{equation}
Its two basic invariants are the \emph{trace form} (a Euclidean, positive-definite bilinear
form) and the \emph{cubic norm}:
\begin{equation}
\label{eq:potential:invariants}
\langle X,Y\rangle=\operatorname{Tr}(X\circ Y),\qquad
Q(X):=\operatorname{Tr}(X^2)=\langle X,X\rangle\ge0,\qquad
N(X)=\det X\ \ (\text{cubic}).
\end{equation}
With the linear trace $\operatorname{tr}(X)=\xi_1+\xi_2+\xi_3$ and identity
$E=\operatorname{diag}(1,1,1)$, every $X$ obeys the cubic characteristic identity
$X^{\circ3}-\operatorname{tr}(X)\,X^{\circ2}+\sigma_2(X)\,X-N(X)\,E=0$ with
$\sigma_2(X)=\tfrac12(\operatorname{tr}(X)^2-Q(X))$. The Jordan eigenvalues
$\lambda_1,\lambda_2,\lambda_3\in\mathbb R$ satisfy
$\operatorname{tr}(X)=\sum_i\lambda_i$, $Q(X)=\sum_i\lambda_i^2$,
$N(X)=\lambda_1\lambda_2\lambda_3$.

The \emph{sharp map} (Jordan adjoint) $X\mapsto X^\#$ is the unique quadratic map with
\begin{equation}
\label{eq:potential:sharp}
X^\#\circ X=N(X)\,E,\qquad (X^\#)^\#=N(X)\,X,\qquad X^\#=\nabla N(X),
\end{equation}
and is the matrix $2\times2$-cofactor (adjugate): in a Jordan frame
$\operatorname{diag}(\lambda_1,\lambda_2,\lambda_3)^\#
=\operatorname{diag}(\lambda_2\lambda_3,\lambda_3\lambda_1,\lambda_1\lambda_2)$.

\paragraph{The two $E_6$'s (a locked distinction).}
The reduced structure group is the non-compact $E_{6(-26)}$ (dim $78$); it preserves $N$ only
up to a scalar character, $N(gX)=\lambda(g)N(X)$, and covariantly transforms the sharp map.
It does \emph{not} preserve $Q=\operatorname{Tr}(X^2)$; indeed $\operatorname{Aut}(J_3(\mathbb O))=F_4$
(compact, dim $52$) is exactly the subgroup preserving both $Q$ and $N$. We use $E_{6(-26)}$
solely to classify orbits of the real $27$ by \textbf{rank and $N$}. The \emph{dynamics}, by
contrast, live on the \textbf{compact} $E_6$: there $H=\mathrm{Spin}(10)\times U(1)$ is
compact, so the coset metric is positive-definite and the gauge sector unitary; on the
compact form $Q=\langle X,X\rangle$ is the invariant trace form and the norm character is
identically $1$. Orbits are classified by rank and $N$ (\emph{not} by $(Q,N)$):
\begin{equation}
\label{eq:potential:rank}
\operatorname{rank}(X)\le1\iff X^\#=0,\qquad
\operatorname{rank}(X)\le2\iff N(X)=0,\qquad
\operatorname{rank}(X)=3\iff N(X)\neq0.
\end{equation}

\subsection{The potential class (Postulate P1)}
\label{sec:potential:class}

\begin{postulate}[Potential class, P1]
\label{post:potential:P1}
The order parameter $\Phi$ is governed by the $E_6$-invariant polynomial potential
\begin{equation}
\label{eq:potential:v}
V(X)=\alpha\,\langle X^\#,X^\#\rangle
+\beta\,\bigl(\operatorname{Tr}(X^2)-c_2\bigr)^2
+\gamma\,\bigl(\operatorname{Tr}(X^2)\bigr)^k,
\qquad \alpha,\beta,\gamma>0,\ \ k\ge2\ \text{even},
\end{equation}
built from the cubic norm (through the sharp map $X^\#=\nabla N$) and the trace form $Q$.
On the compact real form every term is a genuine $E_6$-invariant.
\end{postulate}

The first term is genuinely $E_6$-covariant and vanishes precisely on rank $\le1$; the
$Q$-terms are $F_4$-invariant selectors that fix the scale and the representative. Postulate
P1 asserts that one is \emph{allowed} to write such a potential; the theorems below establish
\emph{what it selects}.

\subsection{The vacuum-selection theorem (P1 \texorpdfstring{$\Rightarrow$}{=>} P2)}
\label{sec:potential:theorem}

\begin{lemma}[Zero locus of the sharp map]
\label{lem:potential:L1}
For $X\in J_3(\mathbb O)$, \ $X^\#=0\iff\operatorname{rank}(X)\le1$. The normalized
($\operatorname{tr}=1$) rank-$1$ elements are exactly the primitive idempotents, forming the
single $F_4$-orbit $F_4/\mathrm{Spin}(9)\cong\mathbb{OP}^2$ (real dim $16$); the
projective $E_6$-orbit of a rank-$1$ line in $\mathbb P(27_{\mathbb C})$ is the complex
Cayley plane $\mathrm{EIII}=E_6/(\mathrm{Spin}(10)\times U(1))$, real dim $32$.
\end{lemma}

\begin{proof}
In a Jordan frame $X=\operatorname{diag}(\lambda_1,\lambda_2,\lambda_3)$, the sharp map gives
$X^\#=\operatorname{diag}(\lambda_2\lambda_3,\lambda_3\lambda_1,\lambda_1\lambda_2)$, which
vanishes iff at most one $\lambda_i\neq0$, i.e. $\operatorname{rank}(X)\le1$; $F_4$-equivariance
transports this to all of $J_3(\mathbb O)$. The orbit statements are the standard
identifications of the real rank-$1$ idempotents ($\mathbb{OP}^2$, dim $16$) and the minimal
$E_6$-orbit / complex Cayley plane (EIII, dim $32$).
\end{proof}

\begin{theorem}[Positivity of the sharp term]
\label{thm:potential:T1}
For all $X\in J_3(\mathbb O)$, \ $\langle X^\#,X^\#\rangle\ge0$, with equality iff $X^\#=0$
iff $\operatorname{rank}(X)\le1$. Explicitly
\begin{equation}
\label{eq:potential:sharpsq}
\langle X^\#,X^\#\rangle=(\lambda_2\lambda_3)^2+(\lambda_3\lambda_1)^2+(\lambda_1\lambda_2)^2
=\sigma_2(X)^2-2\,\operatorname{tr}(X)\,N(X)\ \ (\text{a sum of squares}).
\end{equation}
\end{theorem}

\begin{proof}
The trace form is positive-definite on the Euclidean Jordan algebra, so $\langle Y,Y\rangle\ge0$
with equality iff $Y=0$; take $Y=X^\#$ and apply \autoref{lem:potential:L1}. The frame
expression \eqref{eq:potential:sharpsq} is manifestly a sum of squares.
\end{proof}

\begin{theorem}[Compact-$E_6$ invariance of the sharp term]
\label{thm:potential:T2}
The scalar $\langle X^\#,X^\#\rangle$ is invariant under the compact real form $E_6$ (and
$E_{6(-14)}$): \ $\langle(gX)^\#,(gX)^\#\rangle=\langle X^\#,X^\#\rangle$ for all
$g\in E_6^{\,\mathrm{compact}}$.
\end{theorem}

\begin{proof}
On the compact form the $27$ is unitary, with invariant positive-definite form whose real
part is the trace form; hence the compact $E_6$ preserves $\langle\cdot,\cdot\rangle$. As a
degree-$4$ relative invariant, $\langle X^\#,X^\#\rangle$ transforms by a power of the norm
character $\lambda(g)$; but a continuous homomorphism from a compact connected group to
$\mathbb R_{>0}$ is identically $1$, so $\lambda\equiv1$ and the relative invariant is a
genuine invariant. (Under $E_{6(-26)}$, $\lambda$ is unbounded and $\langle\cdot,\cdot\rangle$
is not preserved; this is precisely why the dynamics use the compact form.)
\end{proof}

\paragraph{Status of $Q$.}
$Q=\operatorname{Tr}(X^2)$ is $F_4$-invariant only and is \emph{not} preserved by the
non-compact structure group; we never claim otherwise, nor that $(Q,N)$ separates $E_6$-orbits.
On the compact form, however, $Q=\langle X,X\rangle$ \emph{is} the invariant trace form, so
$V$ in \eqref{eq:potential:v} is a genuine compact-$E_6$ invariant potential.

\begin{lemma}[The selector has a unique positive minimizer]
\label{lem:potential:L2}
Set $W(q)=\beta(q-c_2)^2+\gamma q^k$ for $q\ge0$, $\beta,\gamma>0$, $k\ge2$ even. There is a
$c_2^{\,\star}>0$ such that for $c_2>c_2^{\,\star}$, $W$ is strictly convex on $[0,\infty)$ and
attains its global minimum at a unique interior point $q_\star\in(0,c_2)$ with $W'(q_\star)=0$,
$W''(q_\star)>0$, and $W(q_\star)<W(0)=\beta c_2^2$. Moreover
\begin{equation}
\label{eq:potential:qstar}
q_\star=c_2-\frac{\gamma k}{2\beta}\,c_2^{\,k-1}+O(c_2^{\,2k-3}),
\qquad q_\star\xrightarrow[\gamma\to0^+]{}c_2.
\end{equation}
\end{lemma}

\begin{proof}
$W(0)=\beta c_2^2>0$ and $W(q)\to+\infty$, so a minimum exists on a compact set. Since
$W''(q)=2\beta+\gamma k(k-1)q^{k-2}>0$ for all $q\ge0$ ($k\ge2$ even), $W$ is strictly convex;
with $W'(0)=-2\beta c_2<0$ it has a unique stationary point $q_\star>0$, the global minimizer,
and $W(q_\star)<W(0)$. As $W'(c_2)=\gamma k c_2^{k-1}>0$, the root lies below $c_2$; the
expansion follows by solving $2\beta(q_\star-c_2)=-\gamma k q_\star^{k-1}$ perturbatively.
\end{proof}

The selector thus \emph{fixes the scale}: among rank-$1$ elements $X=\lambda P$ (with
$Q=\lambda^2$, $N=0$, $X^\#=0$), the minimizers have $\lambda^2=q_\star$, a fixed nonzero
norm; in particular the origin ($q=0$, $W(0)=\beta c_2^2>W(q_\star)$) is excluded.

\begin{lemma}[Coercivity requires the compact form]
\label{lem:potential:L3}
On $J_3(\mathbb O)$ with the compact-$E_6$ Euclidean norm $\|X\|^2=Q(X)$, the potential $V$ is
coercive and bounded below: $V(X)\ge\gamma\|X\|^{2k}\to+\infty$. On a non-compact real-form
orbit, where $Q$ is not preserved and not a proper exhaustion, $V$ need not be coercive.
\end{lemma}

\begin{proof}
$V\ge\gamma Q^k=\gamma\|X\|^{2k}$ from \eqref{eq:potential:v} ($\alpha$-term and the square
are $\ge0$), giving coercivity on the compact form. Under $E_{6(-26)}$ the character
$\lambda(g)$ is unbounded, so a fixed $X$ with $N(X)\neq0$ can be driven to $Q\to0$ along a
fixed $N$-level set; $Q$ is not proper there and sublevel sets are unbounded.
\end{proof}

\begin{theorem}[Global-minimum locus is the rank-$1$ EIII orbit; P2]
\label{thm:potential:T4}
\label{thm:vacuum:T4}
\label{thm:potential:vacuum}
\label{thm:ExplicitBreaking}
With $\alpha,\beta,\gamma>0$, $k\ge2$ even, and $c_2>c_2^{\,\star}$
(\autoref{lem:potential:L2}), every global minimizer $X_\star$ of $V$ on $J_3(\mathbb O)$
satisfies $X_\star^\#=0$ and $\operatorname{Tr}(X_\star^2)=q_\star$; equivalently
\begin{equation}
\label{eq:potential:argmin}
\arg\min_{X\in J_3(\mathbb O)}V(X)
=\{\,X=\lambda P:\ P\ \text{primitive idempotent},\ \lambda^2=q_\star\,\},
\end{equation}
whose projective $E_6$-orbit is the vacuum manifold
\begin{equation}
\label{eq:potential:vacuum}
\mathcal M=\mathrm{EIII}=\frac{E_6}{\mathrm{Spin}(10)\times U(1)},
\qquad \dim_{\mathbb R}\mathcal M=32.
\end{equation}
This proves postulate~\textbf{P2} from \autoref{post:potential:P1}.
\end{theorem}

\begin{proof}
By \autoref{lem:potential:L3} a global minimum is attained. For any $X$,
\begin{equation*}
V(X)=\underbrace{\alpha\,\langle X^\#,X^\#\rangle}_{\ge0}+W(Q(X))\ge W(Q(X))\ge W(q_\star),
\end{equation*}
using \autoref{thm:potential:T1} and \autoref{lem:potential:L2}. Equality in the first step
forces $X^\#=0$, i.e. $\operatorname{rank}(X)\le1$ (\autoref{lem:potential:L1}); equality in
the second forces $Q(X)=q_\star>0$, excluding $X=0$ and forcing rank exactly $1$. Both hold
iff $X=\lambda P$ with $\lambda^2=q_\star$, and every such $X$ achieves $V=W(q_\star)$. The
projective orbit of these rank-$1$ idempotents is EIII (\autoref{lem:potential:L1}).
\end{proof}

\paragraph{Why rank $2$ and rank $3$ strictly lose.}
For $\operatorname{rank}(X)\ge2$, \autoref{lem:potential:L1} gives $X^\#\neq0$, so
$\langle X^\#,X^\#\rangle>0$ strictly (\autoref{thm:potential:T1}) and
$V(X)>W(q_\star)$ by a deficit at least $\alpha\langle X^\#,X^\#\rangle>0$. In particular
$\operatorname{diag}(1,1,1)$ (rank $3$, $\langle X^\#,X^\#\rangle=3>0$, orbit $E_6/F_4$ of
real dimension $26$) is \emph{not} a minimizer. The vacuum is the rank-$1$ idempotent EIII
orbit (dim $32$), \emph{not} the $E_6/F_4$ orbit of $\operatorname{diag}(1,1,1)$ (dim $26$).

\begin{remark}[Coset tangent and metric]
\label{rmk:potential:tangent}
The tangent space to $\mathcal M$ at the base idempotent is the coset module
\begin{equation}
\label{eq:potential:cosettangent}
\mathfrak m=16_{-3}\oplus\overline{16}_{+3}\quad(32\ \text{real}),
\end{equation}
read off from the adjoint branching $78=45_0\oplus1_0\oplus16_{-3}\oplus\overline{16}_{+3}$
under $H=\mathrm{Spin}(10)\times U(1)$. It is real-irreducible of complex type (commutant
$\mathbb C$) and contains \emph{no} $\mathrm{SO}(10)$ singlet; the $U(1)$ charge forbids the
symmetric $16\times16$ and $\overline{16}\times\overline{16}$ invariants, leaving a unique
charge-$0$ cross-pairing. The induced coset metric $K|_{\mathfrak m}$ is therefore unique and
positive-definite (Schur), and serves as the kinetic/target-space metric for the $32$ clock
fields (Goldstones). It carries \emph{no} Lorentzian content: the $(1,3)$ signature is added
structure (P3), soldered from the cubic norm via $H_2(\mathbb O)\cong\mathbb R^{1,9}$, not a
Killing-form consequence; see \autoref{sec:spacetime:main}.
\end{remark}

\subsection{The symmetric origin is a strict local maximum}
\label{sec:potential:origin}

\begin{theorem}[Origin is a strict local maximum]
\label{thm:potential:T5}
At $X=0$,
\begin{equation}
\label{eq:potential:hess0}
\nabla V(0)=0,\qquad
\operatorname{Hess}V(0)=-4\beta c_2\,\mathbb 1\ \prec\ 0
\quad(\text{negative-definite, }\beta,c_2>0),
\end{equation}
so $X=0$ is a strict local maximum, all $27$ Hessian eigenvalues equal to $-4\beta c_2<0$.
Since $V(0)=\beta c_2^2>W(q_\star)$, the rank-$1$ orbit beats the origin locally and globally.
\end{theorem}

\begin{proof}
$V$ has no linear part. The $\alpha\langle X^\#,X^\#\rangle$ term is homogeneous of degree $4$
and the $\gamma Q^k$ term of degree $2k\ge4$, so both have vanishing Hessian at $0$. Only the
quadratic piece $-2\beta c_2 Q$ of $\beta(Q-c_2)^2$ contributes; with $Q$ having Hessian
$2\,\mathbb 1$ in the trace metric this gives $\operatorname{Hess}V(0)=-4\beta c_2\,\mathbb 1$,
negative-definite.
\end{proof}

\begin{theorem}[Roll-down to $\mathcal M$]
\label{thm:potential:T6}
Under gradient flow $\dot X=-\nabla V(X)$ (or the EOM with damping), the origin is linearly
unstable with growth rate $4\beta c_2>0$ in all $27$ directions; $V$ is a strict Lyapunov
function, coercivity (\autoref{lem:potential:L3}) confines trajectories to compact sublevel
sets, and by LaSalle the $\omega$-limit set lies in $\{\nabla V=0\}$. The unique stable
critical set is the global minimum $\mathcal M$ (\autoref{thm:potential:T4}); the origin
(maximum, \autoref{thm:potential:T5}) and the rank-$2$/rank-$3$ saddles ($X^\#\neq0$) are all
unstable. Hence generic initial data roll down to $\mathcal M=\mathrm{EIII}$.
\end{theorem}

\begin{proof}
Linearization at $0$ uses $\operatorname{Hess}V(0)\prec0$ (\autoref{thm:potential:T5}), giving
positive growth in every direction. $\dot V=-\|\nabla V\|^2\le0$ vanishes only at critical
points; LaSalle plus coercivity confine $\omega$-limits to the critical set, of which only
$\mathcal M$ is Lyapunov-stable. The basin of $\mathcal M$ is open and dense.
\end{proof}

\paragraph{Tree-level Goldstone consistency.}
On $\mathcal M$, $V$ is $E_6$-invariant and constant along the orbit, so
$\operatorname{Hess}V|_{\mathfrak m}=0$ at tree level: the $32$ clock fields are exact
Goldstones. Radiative (Coleman--Weinberg) corrections are $H$-invariant; by Schur on the
real-irreducible module $\mathfrak m$ they give \emph{one} common mass on all $32$ directions.
Of these, $28$ are physical massive scalars (single common mass) and $4$ are the
gauge/soldering modes eaten by diffeomorphism and local Lorentz; this $28+4$ split and the
common radiative mass are carried into the spectrum and induced-gravity analysis of
\autoref{sec:gravity:main}. The origin's $27$ negative modes and the vacuum's flat coset
directions are not in tension: they live at different points of field space.

\subsection{Summary of the logical chain}
\label{sec:potential:summary}

\begin{center}
\begin{tabular}{lll}
\toprule
step & statement & type \\
\midrule
\autoref{lem:potential:L1} & $X^\#=0\iff\operatorname{rank}(X)\le1$; rank-$1$ orbit $=$ EIII (dim $32$) & Theorem \\
\autoref{thm:potential:T1} & $\langle X^\#,X^\#\rangle\ge0$, $=0\iff\operatorname{rank}\le1$ & Theorem \\
\autoref{thm:potential:T2} & $\langle X^\#,X^\#\rangle$ is compact-$E_6$ invariant & Theorem \\
\autoref{lem:potential:L2} & selector $W$ strictly convex, unique $q_\star>0$, $W(q_\star)<W(0)$ & Theorem \\
\autoref{lem:potential:L3} & coercivity holds on the compact form, fails on non-compact orbits & Theorem \\
\autoref{thm:potential:T4} & \textbf{global-min locus $=$ rank-$1$ EIII orbit, dim $32$ (P2)} & \textbf{Theorem} \\
\autoref{thm:potential:T5} & $\operatorname{Hess}V(0)=-4\beta c_2\,\mathbb1\prec0$: origin strict max & Theorem \\
\autoref{thm:potential:T6} & roll-down: generic flows attract to $\mathcal M$ & Theorem \\
\autoref{post:potential:P1} & the potential class is added structure & Postulate \\
\bottomrule
\end{tabular}
\end{center}

Given the postulated potential class P1, \autoref{thm:potential:T4} proves P2: the unique
global-minimum locus of $V$ is the rank-$1$ primitive-idempotent orbit
$\mathcal M=\mathrm{EIII}=E_6/(\mathrm{Spin}(10)\times U(1))$, real dimension $32$. The sharp
term enforces rank $\le1$; the $Q$-selector fixes the nonzero scale and excludes the origin;
the compact real form supplies coercivity and a definite invariant metric; and the origin is
a strict local maximum driving the roll-down. No rank-$2$/rank-$3$ point --- in particular
none on the $E_6/F_4$ orbit of $\operatorname{diag}(1,1,1)$ (dim $26$) --- can be a minimizer.
